%%
% 摘要信息
% 本文档中前缀"c-"代表中文版字段, 前缀"e-"代表英文版字段
% 摘要内容应概括地反映出本论文的主要内容，主要说明本论文的研究目的、内容、方法、成果和结论。要突出本论文的创造性成果或新见解，不要与引言相 混淆。语言力求精练、准确，以 300—500 字为宜。
% 在摘要的下方另起一行，注明本文的关键词（3—5 个）。关键词是供检索用的主题词条，应采用能覆盖论文主要内容的通用技术词条(参照相应的技术术语 标准)。按词条的外延层次排列（外延大的排在前面）。摘要与关键词应在同一页。
% modifier: 黄俊杰(huangjj27, 349373001dc@gmail.com)
% update date: 2017-04-15
%%

\cabstract{
在学习王老师的大数据原理与技术课程\cite{lecture01}后，在此报告中完成相应的课后作业，本作业围绕“图像边缘检测与分类”展开，要求使用 Python 与 OpenCV 完成\textbf{两部分任务}。\textbf{图像边缘检测}：任选一张图像，采用Sobel算法进行边缘检测，并展示原始图像与处理后的分割结果。
\textbf{图像分类对比实验}：在可选数据集（如MNIST或CIFAR-10）上，分别实现一种传统机器学习方法（如支持向量机SVM、决策树等）与一种深度学习方法（如多层感知机MLP、卷积神经网络CNN等），并对两类方法的分类准确率与计算效率进行比较分析。

作业提交要求：
\begin{itemize}
    \item 提交内容：代码与作业报告
    \item 命名格式：“学号+姓名+作业1.zip”
    \item 提交至邮箱：BDSysu2026Spring@163.com
    \item 截止时间：布置作业之后的下一个周日24:00
    \item 严禁抄袭：任何一次抄袭，平时成绩按零分处理
\end{itemize}

\textbf{边缘检测}采用 OpenCV 实现的 Sobel 算子，对示例图像（位于 code/data/test.png）进行了灰度化、平滑与 Sobel 卷积操作，得到边缘强度图并保存为 code/data/test\_sobel.png。图像分类部分在 MNIST 数据集上进行比较实验：

\textbf{分类对比实验}以两类代表性数据集为目标：一是 Digits（小样本、低维），二是 MNIST（大样本、高维）；对比对象为传统机器学习方法（SVM）与深度学习方法（MLP）。主要实验结果如下：

（1）Digits：SVM 准确率 0.9889，训练耗时约 0.06s，推理约 0.02s；MLP（小型网络）准确率 0.9694，训练耗时较短，推理吞吐较高；

（2）MNIST：在完整 MNIST 上 MLP 准确率 0.9611，训练耗时显著增加（约 33.6s）；在未经专门调参的 MNIST 子集上直接使用 SVM 表现较差（准确率约 0.1135）且计算代价增加。

实验表明：在小样本、低维任务中，传统方法（如 SVM）在准确率与计算效率上具有优势；而在大规模、高维任务中，深度学习模型更能发挥表达与扩展能力，但代价是更高的训练开销。
}
% 中文关键词(每个关键词之间用“，”分开,最后一个关键词不打标点符号。)
\ckeywords{图像边缘检测，Sobel算子，图像分类，SVM与CNN，MNIST与CIFAR-10}
% 英文关键词(每个关键词之间用,分开, 最后一个关键词不打标点符号。)
\ekeywords{}


